\documentclass{article}
\usepackage{amsmath,amssymb}
\usepackage{kotex}
\begin{document}

\noindent\textbf{문제.} 응범이는 폭풍우 속에서 공정한 동전을 반복해서 던지고 있다. 동전을 한 번 던질 때마다 \(4\%\)의 확률로 동전이 바람에 날아가 버린다. 응범이는 동전을 잃어버리기 전까지 나온 앞면의 횟수를 \(H\), 뒷면의 횟수를 \(T\)로 기록한다. 단, 동전을 던지는 순간 동전이 바람에 날아간 경우에는 그 시행의 결과를 기록하지 않는다. 이때
\[ |H-T| \]
의 기댓값을 구하여라.

\noindent\textbf{해설.}
첫 번 째 방법

동전이 한 번의 시행에서 바람에 날아갈 확률을
\[ p=\frac{1}{25} \]
라 하자. 또한
\[ D=|H-T| \]
라 하자.

동전이 바람에 날아가지 않고 정상적으로 던져졌다고 하자. 현재 \(D\neq 0\)이면, 다음 동전의 결과에 따라 \(D\)는 같은 확률로 \(1\) 증가하거나 \(1\) 감소한다. 따라서 한 번의 동전 던지기로 인한 \(D\)의 변화량의 기댓값은 \(0\)이다.

반면 \(D=0\)인 상태에서 동전을 정상적으로 한 번 던지면 반드시 \(D=1\)이 되므로, 이 경우에는 \(D\)가 \(1\)만큼 증가한다. 따라서 최종적인 \(D\)의 기댓값은 \(D=0\)인 상태에서 실제로 동전을 한 번 더 던지게 되는 횟수의 기댓값과 같다.

\(D=0\)인 상태는 앞면과 뒷면의 수가 같을 때이므로, 정상적인 동전을 \(2k\)번 던진 뒤에만 일어날 수 있다. 처음 \(2k\)번의 시행에서 동전이 날아가지 않을 확률은
\[ (1-p)^{2k} \]
이고, 그 \(2k\)번의 동전 던지기에서 앞면과 뒷면이 각각 \(k\)번씩 나올 확률은
\[ \frac{\binom{2k}{k}}{2^{2k}}. \]
또한 그다음 시행에서도 동전이 날아가지 않아야 실제 동전 던지기가 한 번 더 이루어지므로, 그 확률은 \(1-p\)이다. 따라서
\[
\begin{aligned}
\mathbb{E}(|H-T|)
&=(1-p)\sum_{k=0}^{\infty} (1-p)^{2k}\frac{\binom{2k}{k}}{2^{2k}}\\
&=(1-p)\sum_{k=0}^{\infty} \binom{2k}{k} \left(\frac{(1-p)^2}{4}\right)^k.
\end{aligned}
\]
여기서 생성함수
\[ \sum_{k=0}^{\infty}\binom{2k}{k}x^k =\frac{1}{\sqrt{1-4x}} \]
를 이용하면
\[
\begin{aligned}
\mathbb{E}(|H-T|)
&=\frac{1-p}{\sqrt{1-(1-p)^2}}\\
&=\frac{1-p}{\sqrt{p(2-p)}}.
\end{aligned}
\]
이제 \(p=\frac{1}{25}\)를 대입하면
\[
\begin{aligned}
\mathbb{E}(|H-T|)
&=\frac{\frac{24}{25}} {\sqrt{\frac{1}{25}\cdot\frac{49}{25}}}\\
&=\frac{24}{7}.
\end{aligned}
\]
따라서 구하는 기댓값은
\[ \boxed{\frac{24}{7}} \]
이다.

두 번째 방법

각 정수 \(n\)에 대하여
\[ E_n=\mathbb{E}(|H-T+n|) \]
이라 하자. 대칭성에 의해
\[ E_{-n}=E_n \]
이다.

현재 \(H-T=n\)인 상태에서 다음 시행을 생각하자. 동전이 바람에 날아갈 확률은 \(p\)이고, 날아가지 않을 확률은 \(1-p\)이다. 동전이 날아가지 않은 경우에는 공정한 동전을 던지므로 \(H-T\)가 \(1\) 증가하거나 \(1\) 감소할 확률이 각각 \(\frac12\)이다.

따라서 모든 \(n\ge 0\)에 대하여
\[ 2E_n=(1-p)E_{n-1}+(1-p)E_{n+1}+2pn \]
이 성립한다.

이제
\[ \alpha=\frac{1-\sqrt{p(2-p)}}{1-p} \]
라 하자. 이는
\[ (1-p)x^2-2x+(1-p)=0 \]
의 두 근 중 \(1\)보다 작은 근이다.

위의 점화식에 \(\alpha^n\)을 곱하여 \(n=1\)부터 합하면
\[ \sum_{n=1}^{\infty}2\alpha^nE_n = \sum_{n=1}^{\infty} \alpha^n\bigl((1-p)E_{n-1}+(1-p)E_{n+1}+2pn\bigr). \]
우변을 정리하면
\[
\begin{aligned}
\sum_{n=1}^{\infty}2\alpha^nE_n
&=\alpha(1-p)E_0+\alpha^2(1-p)E_1 +\sum_{n=1}^{\infty}2pn\alpha^n\\
&\quad +\sum_{n=2}^{\infty} \bigl((1-p)\alpha^{n-1}+(1-p)\alpha^{n+1}\bigr)E_n.
\end{aligned}
\]
그런데 \(\alpha\)는
\[ (1-p)\alpha^2-2\alpha+(1-p)=0 \]
을 만족하므로
\[ (1-p)\alpha^{n-1}+(1-p)\alpha^{n+1} =2\alpha^n \]
이다. 따라서 \(n\ge2\)인 항들은 양변에서 소거되어
\[ (1-p)E_1-\alpha(1-p)E_0 = \sum_{n=1}^{\infty}2pn\alpha^n. \]
또한
\[ \sum_{n=1}^{\infty}n\alpha^n = \frac{\alpha}{(1-\alpha)^2} \]
이므로
\[ (1-p)E_1-\alpha(1-p)E_0 = \frac{2p\alpha}{(1-\alpha)^2}. \]
한편 \(n=0\)일 때 점화식은
\[ 2E_0=(1-p)E_{-1}+(1-p)E_1 \]
이고 \(E_{-1}=E_1\)이므로
\[ E_0=(1-p)E_1. \]
이를 위 식에 대입하면
\[ E_0\bigl(1-\alpha(1-p)\bigr) = \frac{2p\alpha}{(1-\alpha)^2}. \]
이제
\[ p=\frac1{25} \]
을 대입하면
\[ \alpha = \frac{1-\sqrt{\frac1{25}\cdot\frac{49}{25}}} {\frac{24}{25}} = \frac34. \]
따라서
\[ E_0 \left( 1-\frac34\cdot\frac{24}{25} \right) = \frac{2\cdot\frac1{25}\cdot\frac34} {\left(1-\frac34\right)^2}. \]
이를 계산하면
\[ E_0=\frac{24}{7}. \]
처음에는 \(H=T=0\)이므로 구하는 기댓값은 \(E_0\)이다. 따라서
\[ \boxed{\mathbb{E}(|H-T|)=\frac{24}{7}}. \]

\end{document}
